
\section{parameters.xml}

\subsection{RBAMetabolism container}

The metabolism file is strongly inspired by SBML.\@
More precisely, it is a subpart of an SBML file.

\subsubsection{RBAMetabolism}
\label{sec:rba_metabolism}

The outermost portion of the metabolism file is an instance of class
\rbametabolism, shown in Figure~\ref{fig:metabolism_doc}.

\begin{figure}
  \centering
  \includegraphics[scale=0.9]{figures/metabolism_doc}
  \caption{XML structure of metabolism document.}
\label{fig:metabolism_doc}
\end{figure}

Currently, \rbametabolism{} has no simple attributes.
It includes exactly one instance of \textbf{ListOf} container classes.
All \textbf{ListOf} classes do not have own attributes,
they are merely used to organize a list of instances from another class.
This organization was inspired by SBML.\@

\subsubsection{Compartment}
\label{sec:compartment}

The \compartment{} class is used to list existing cell compartments.

\paragraph{The \textit{id} attribute}
The \textbf{id} attribute is a string defining the identifier of a compartment.

\subsubsection{Species}
\label{sec:species}

The \species{} class is used to list \emph{metabolic} species.

\paragraph{The \textit{id} attribute}
The \textbf{id} attribute is a string defining the identifier of a metabolite.

\paragraph{The \textit{boundaryCondition} attribute}
The \textbf{boundaryCondition} attribute is a boolean.
If the attribute is set to true, the metabolite is considered to be at
a constant concentration.
In other words, it is not affected by reactions.
This is typical for metabolites in the external medium.
